% Ad Familiares 15.16 --- headnote
% ad-familiares-15-16, January 45 BC, Rome

\textit{Cicero to C. Cassius Longinus, written from
Rome before the kalends of January 709 AUC ---
month-precision January 45 BC, the Perseus dateline
giving only \textit{ante m.\ Ian.} The first of the
three letters in this cluster in which Cicero engages
Cassius on his recent conversion to Epicureanism. The
joke runs through the opening: Cicero, writing to a
silent correspondent, imagines Cassius vividly present
as he writes --- and disclaims that this presence
arrives by Democritean \textit{[Greek: eidōla]}, the
atomic images of the school, here teased through the
Latin coinage \textit{spectra} of Catius the Insubrian,
the recently dead Epicurean popularizer.}

\textit{The second section presses the joke
philosophically: even if the eye could be struck by
\textit{spectra}, the mind cannot. Cicero asks whether
Cassius's \textit{spectrum} is at his command, and
whether the island of Britain --- a place neither has
seen but both can think of --- will send him its
\textit{[Greek: eidōlon]} at will. The third turns into
the running gag of the diptych: a mock-interdict
demanding Cassius's restoration to the school
\textit{[Greek: hairesei]} from which the men under
arms (Caesar's victory) have ejected him, written in
the language of the praetor's interdict ``vi hominibus
armatis.'' The letter ends in the period's standard
guarded coda: he writes about philosophy because the
state cannot be written about.}
